\documentclass[10pt,a4paper]{amsart}
\usepackage[margin=19mm]{geometry}
\usepackage{amsmath,amssymb,mathtools}
\usepackage[scr=rsfs,cal=euler]{mathalfa}
\usepackage{xfrac}
\usepackage[unicode,hidelinks,bookmarks=false]{hyperref}
\newcommand{\ie}{i.e.\ }
\newcommand{\cf}{cf.\ }
\newcommand{\resp}{respectively\ }
\newcommand{\Subsection}{\subsection}
\providecommand{\nobreakdash}{\nobreak\hbox{-}}
\setlength{\emergencystretch}{2em}
\multlinegap0pt
\usepackage{amssymb}
\usepackage{mathtools}
\usepackage{tikz}
\usepackage{xcolor}
\newtheorem{assumption}{Assumption}
\newtheorem{definition}{Definition}
\newtheorem{lemma}{Lemma}
\newtheorem{proposition}{Proposition}
\newcommand{\Hcauchy}{\mathcal H_{\mathrm{C},N}}
\newcommand{\Hscatt}{\mathcal H_{\mathrm{sc,deg}}}
\newcommand{\Lin}{\mathsf{Lin}}
\newcommand{\ME}{\mathsf{ME}}
\newcommand{\dd}{\mathrm{d}}
\title{\Large\textbf{Coulomb Sectors and Scattering for Maxwell-Higgs Fields on Schwarzschild and Slowly Rotating Kerr Backgrounds}}
\author{Bobby Eka Gunara, Mulyanto, Fiki Taufik Akbar}
\date{Source: arXiv:2603.04202v2 (CC BY 4.0)}
\begin{document}
\maketitle

\noindent\emph{Source and licence.} \Large\textbf{Coulomb Sectors and Scattering for Maxwell-Higgs Fields on Schwarzschild and Slowly Rotating Kerr Backgrounds}. Version arXiv:2603.04202v2 (CC BY 4.0), distributed by arXiv under the Creative Commons Attribution 4.0 International licence, http://creativecommons.org/licenses/by/4.0/, as recorded in the arXiv licence metadata for this exact version and shown on the abstract page https://arxiv.org/abs/2603.04202v2. Full licence notice: \texttt{licenses/arxiv-2603.04202-CC-BY-4.0.txt}.\par\smallskip

\noindent\textit{Editorial scope.} Counted unit: the article's proposition prop:slow-core-estimates-full (source0.tex lines 4267--4305). The article's complete original proof of it is delivered with it (source0.tex lines 4307--4418, one \texttt{proof} environment of 112 lines). No mathematics is written, repaired or re-derived: the statement and the proof are the article's own lines, in the article's own notation, and no retained line is substituted or re-flowed.  Retained uncounted as the source-local context the proof needs: lines 509--534; lines 609--645; lines 653--657; lines 913--935; lines 998--1000; lines 1033--1107; lines 1595--1606; lines 3931--4048; lines 4050--4082; lines 4088--4099; lines 4171--4175; lines 4234--4241.  Nothing was omitted from any retained span, so the package carries no omission record.  No retained line contains a citation, so the wrapper carries no bibliography and no bibliography entry.  No retained line contains a figure environment, an image-inclusion command or drawing code, so no image is delivered and the package declares no asset dependency.  Editorial formatting only, in addition: an AMS \texttt{amsart} preamble that restates the article's own declarations and macros used by the retained lines, namely the environments \texttt{assumption}, \texttt{definition}, \texttt{lemma}, \texttt{proposition} on separate counters, together with the standard packages those lines need.  The retained environments print in this document as Assumption 1, Definition 1, Lemma 1, Proposition 1 in order of appearance, whereas the article prints the same items with the numbering of its own full document.  The complete native source of the article as distributed in the arXiv e-print for this exact version is bundled unchanged as \texttt{sources/195723/source0.tex}.
\medskip
\noindent\textbf{Derivatives, gauge-covariant fields, and currents.}
We use the Lorentzian signature $(-,+,+,+)$ and the Einstein summation convention.
We denote by $\nabla$ the Levi-Civita connection and by $\square_{g}=\nabla^{\mu}\nabla_{\mu}$ the wave operator.
Given a potential $A$ we write $F=\dd A$ for its curvature and $D_{\mu}=\nabla_{\mu}-iA_{\mu}$ for the gauge-covariant derivative.
We reserve the letter $F$ for the curvature $2$-form; when discussing inhomogeneous scalar wave/Klein-Gordon equations we denote the forcing term by $\mathfrak{F}$.
If $\Sigma$ is a spacelike hypersurface with future unit normal $n_{\Sigma}$, then for any current $J$ we write
$\int_{\Sigma}J\cdot n_{\Sigma}$ as shorthand for $\int_{\Sigma}J_{\mu}n_{\Sigma}^{\mu}\,\dd\mu_{\Sigma}$, where $\dd\mu_{\Sigma}$ is the induced volume form.
The Maxwell-Higgs current is
\begin{equation}\label{eq:mh-current}
J_{\mu}:=-i\bigl(D_{\mu}\phi\,\overline\phi-\phi\,\overline{D_{\mu}\phi}\bigr)=2\,\Im(\overline\phi\,D_{\mu}\phi).
\end{equation}
The Lagrangian density and Euler-Lagrange equations are normalized by
\begin{equation}\label{lagrangian}
\mathcal L[A,\phi]
=-\frac14 F_{\alpha\beta}F^{\alpha\beta}
-g^{\alpha\beta}D_{\alpha}\phi\,\overline{D_{\beta}\phi}
-P(\phi,\overline\phi).
\end{equation}
Thus the geometric Maxwell-Higgs system is
\begin{eqnarray}
\nabla^{\mu}F_{\mu\nu}&=&J_{\nu},
\label{eom1}\\
D^{\mu}D_{\mu}\phi&=&\partial_{\overline\phi}P(\phi,\overline\phi).
\label{eom2}
\end{eqnarray}

\begin{assumption}[Scalar potential]\label{asumsiP}
The scalar potential $P:\mathbb{C}\to\mathbb{R}$ is gauge invariant and depends only on the scalar $|\phi|^{2}$, i.e.\ $P(\phi,\bar\phi)=\widetilde P(|\phi|^{2})$ for some $\widetilde P:[0,\infty)\to\mathbb{R}$.
Assume that $\widetilde P\in C^\infty([0,\infty))$, $\widetilde P(0)=0$, and $\widetilde P\ge 0$ (after fixing the irrelevant additive constant so that the potential energy density vanishes in the vacuum configuration).

In addition, we assume that there exist a constant $m^{2}\ge 0$, an integer $N_{P}\ge 0$, and a function $Q\in C^\infty([0,\infty))$ such that
\begin{equation}\label{eq:potential-structure}
\widetilde P'(s)=m^{2}+s\,Q(s)\qquad\text{for all }s\ge 0,
\end{equation}
and, for every integer $j\ge0$, there exists a constant $C_{P,j}$ such that
\begin{equation}\label{eq:Q-growth}
|Q^{(j)}(s)|\le C_{P,j}\,(1+s^{N_{P}})\qquad\text{for all }s\ge 0.
\end{equation}
(In this case, the nonlinear part of the force term $\partial_{\bar\phi}P$ is at least cubic in $\phi$ and has at most polynomial growth in $|\phi|$.)

For concreteness one may keep in mind the following model potentials:
\begin{eqnarray}
\widetilde P(s)
&=&
m^2s+\sum_{n=p}^{N}\alpha_{n}s^{n},
\qquad m^2\ge0,\quad N\ge p\ge2,\quad \alpha_n\ge 0,
\label{V4}
\\
\widetilde P(s)
&=&
c_{3}\bigl(1-\cos(\eta\sqrt{s})\bigr),
\qquad c_3\ge 0,\ \eta>0,
\label{sine}
\\
\widetilde P(s)
&=&
c_{4}\bigl(1-e^{-\lambda s}\bigr),
\qquad c_4\ge 0,\ \lambda>0,
\label{toda}
\end{eqnarray}
which all satisfy \eqref{eq:potential-structure}-\eqref{eq:Q-growth} for suitable choices of $m^{2}\ge0$, $N_{P}$, $Q$, and constants $C_{P,j}$.
(Adding a constant to $\widetilde P$ does not change \eqref{eom2}, and we normalize it so that the potential energy density vanishes in the vacuum.)
\end{assumption}

\begin{equation}\label{eq:potential-remainder}
\partial_{\bar\phi}P(\phi,\bar\phi)=m^{2}\phi+\mathcal R_{P}(\phi),
\qquad
\mathcal R_{P}(\phi):=Q(|\phi|^{2})\,|\phi|^{2}\phi.
\end{equation}

\begin{definition}[Admissible stationary black-hole exterior]\label{def:admissible-exterior}
An \emph{admissible stationary black-hole exterior} consists of an asymptotically flat stationary black-hole exterior $(\mathcal D,g)$ together with an \emph{admissible structure}
\begin{equation}
\bigl(\{\Sigma_{\tau}\}_{\tau\in\mathbb R},\,\mathcal R(\tau_{0},\tau_{1}),\,r,\,N,\,\mathcal Z\bigr)
\end{equation}
satisfying the following properties.
\begin{enumerate}
\item[(G1)] \textbf{Foliation and slabs.}
$\{\Sigma_{\tau}\}$ is a smooth spacelike foliation of $\mathcal D$ by Cauchy hypersurfaces.
For $\tau_{1}\ge\tau_{0}$, the corresponding slab $\mathcal{R}(\tau_{0},\tau_{1})$ has boundary components
$\Sigma_{\tau_{0}}$, $\Sigma_{\tau_{1}}$, and the portions of the event horizons and null infinities
$\mathcal H^{\pm}(\tau_{0},\tau_{1})$ and $\mathcal I^{\pm}(\tau_{0},\tau_{1})$.
\item[(G2)] \textbf{Asymptotic radius and stationarity.}
In the asymptotically flat region there is a smooth radial coordinate $r$ with $r\sim |x|$ and a complete stationary Killing field $T$
which is timelike for $r$ sufficiently large.
\item[(G3)] \textbf{Redshift multiplier.}
There exists a globally timelike future-directed multiplier field $N$ which agrees with $T$ for $r$ sufficiently large
and is a redshift vector field in a neighbourhood of $\mathcal H^{+}$ (and, by time-reversal, of $\mathcal H^{-}$).
\item[(G4)] \textbf{Commutations.}
$\mathcal Z$ is a fixed finite collection of admissible first-order commutation vector fields used to define higher-order energies and to commute the Maxwell field and the nonlinear system. When additional hidden symmetries are used in the scalar linear theory (for instance the Carter operator on Kerr), they are treated as auxiliary scalar commutators and are not included in $\mathcal Z$.
\end{enumerate}
Whenever we speak of an admissible exterior we implicitly fix one such admissible structure.
\end{definition}

\begin{definition}[Kerr scattering topology]\label{def:kerr-scattering-topology}
On a slowly rotating Kerr exterior the forward Cauchy problem and the two-sided final-state problem are measured in different but compatible topologies. The forward Cauchy space ${\Hcauchy}^{k}(\Sigma_{\tau_0})$ is the completion of smooth data in the nondegenerate order-$k$ $N$-energy norm. The scattering/final-state Cauchy space $\Hscatt^{k}(\Sigma_{\tau_0})$ is the completion induced by the scattering norm whose horizon component is degenerate in the redshift direction. The Cauchy-to-radiation map is continuous from ${\Hcauchy}^{k}$ to the radiation space, while the inverse final-state map is continuous from radiation data to $\Hscatt^{k}$. A final-state trace in ${\Hcauchy}^{k}$ requires an extra redshift-regularity condition on the horizon radiation field. This is the topology used in the zero-sector Kerr instances of $\Lin_k$.
\end{definition}

\begin{definition}[Linear estimates $\Lin_{k}$]\label{def:lin-estimates}
Fix $k\in\mathbb N$.
Let $(\mathcal D,g)$ be an admissible stationary black-hole exterior with admissible structure as in Definition~\ref{def:admissible-exterior}.
We say that $(\mathcal D,g)$ satisfies the \emph{analytic} linear estimates collected in $\Lin_{k}$ if there exist constants $C>0$ and $R>0$ and, for every $\tau_{1}\ge\tau_{0}$, a pair of nonnegative functionals
\begin{equation}
\|\cdot\|_{\mathcal S_{k}(\tau_{0},\tau_{1})},\qquad
\|\cdot\|_{\mathcal X_{k}(\tau_{0},\tau_{1})},
\end{equation}
called the \emph{source norm} and the \emph{solution/bulk norm}, such that the following hold.

\smallskip
\noindent\textbf{Admissibility of the norms.}
The families $\mathcal S_{k}$ and $\mathcal X_{k}$ satisfy the following structural properties.
\begin{enumerate}
\item[(A1)] \textbf{Time monotonicity.}
If $[\tau_{0}',\tau_{1}']\subset[\tau_{0},\tau_{1}]$, then
\begin{equation}
\|\mathfrak{F}\|_{\mathcal S_{k}(\tau_{0}',\tau_{1}')} \le \|\mathfrak{F}\|_{\mathcal S_{k}(\tau_{0},\tau_{1})},
\qquad
\|\psi\|_{\mathcal X_{k}(\tau_{0}',\tau_{1}')} \le \|\psi\|_{\mathcal X_{k}(\tau_{0},\tau_{1})},
\end{equation}
and similarly for Maxwell fields and currents.
\item[(A2)] \textbf{Time localization.}
For $\tau_{0}\le\tau_{1}\le\tau_{2}$ there exists a universal constant $C_{\mathrm{loc}}$ such that
\begin{equation}
\|\mathfrak{F}\|_{\mathcal S_{k}(\tau_{0},\tau_{2})}^{2}
\le
\|\mathfrak{F}\|_{\mathcal S_{k}(\tau_{0},\tau_{1})}^{2}
+
\|\mathfrak{F}\|_{\mathcal S_{k}(\tau_{1},\tau_{2})}^{2},
\end{equation}
and
\begin{equation}
\|\psi\|_{\mathcal X_{k}(\tau_{0},\tau_{2})}^{2}
\le
C_{\mathrm{loc}}\Bigl(
\|\psi\|_{\mathcal X_{k}(\tau_{0},\tau_{1})}^{2}
+
\|\psi\|_{\mathcal X_{k}(\tau_{1},\tau_{2})}^{2}
\Bigr),
\end{equation}
with the analogous property for Maxwell fields.
\item[(A3)] \textbf{Coercivity and radiation fields.}
Finiteness of $\|\psi\|_{\mathcal X_{k}(\tau_{0},\tau_{1})}$ controls the order-$k$ $N$-energy fluxes through $\Sigma_{\tau}$ for $\tau\in[\tau_{0},\tau_{1}]$.
In addition, $\|\psi\|_{\mathcal X_{k}(\tau_{0},\tau_{1})}$ contains a far-field component sufficient to yield existence of order-$k$ radiation fields on $\mathcal I^{\pm}(\tau_{0},\tau_{1})\cup\mathcal H^{\pm}(\tau_{0},\tau_{1})$
together with quantitative control of the corresponding asymptotic energy fluxes.
The same holds for uncharged Maxwell fields.
\end{enumerate}

\begin{enumerate}
\item[(L1)] \textbf{Inhomogeneous energy-Morawetz-$r^{p}$ estimate.}
Let $\psi$ solve the scalar wave equation $\square_{g}\psi=\mathfrak{F}$ on $\mathcal{R}(\tau_{0},\tau_{1})$.
Then
\begin{equation}\label{eq:Lin-abstract}
\|\psi\|_{\mathcal X_{k}(\tau_{0},\tau_{1})}^{2}
\le
C\Bigl(E^{N}_{k}[\psi](\tau_{0})
+\|\mathfrak{F}\|_{\mathcal S_{k}(\tau_{0},\tau_{1})}^{2}\Bigr).
\end{equation}
An analogous estimate holds for (uncharged) Maxwell fields $G$ solving
$\nabla^{\mu}G_{\mu\nu}=J_{\nu}$, $\nabla_{[\alpha}G_{\beta\gamma]}=0$, namely
\begin{equation}
\|G\|_{\mathcal X_{k}(\tau_{0},\tau_{1})}^{2}
\le
C\Bigl(E^{N}_{k}[G](\tau_{0})
+\|J\|_{\mathcal S_{k}(\tau_{0},\tau_{1})}^{2}\Bigr).
\end{equation}

\item[(L2)] \textbf{Linear scattering and final-state well-posedness.}
The Cauchy-to-radiation map extends continuously from the order-$k$ forward Cauchy space ${\Hcauchy}^{k}(\Sigma_{\tau_0})$ to a Hilbert space of linear asymptotic data on $\mathcal I^{\pm}\cup\mathcal H^{\pm}$. The inverse final-state problem on $\mathcal I^{\pm}\cup\mathcal H^{\pm}$ is well posed into the scattering Cauchy space ${\Hscatt}^{k}(\Sigma_{\tau_0})$, with estimates compatible with (L1). In geometries without a horizon blue-shift obstruction these two spaces may coincide; on Kerr they are kept distinct as in Definition~\ref{def:kerr-scattering-topology}.
\end{enumerate}

All Maxwell estimates in $\Lin_{k}$ are understood for the \emph{radiative} (uncharged) part of the field:
for general data with nonzero electric charge one first subtracts the corresponding stationary Coulomb solution.
\end{definition}

\begin{lemma}[Tame commutator bounds for the potential force]\label{lem:potential-tame}
Assume that $P$ satisfies Assumption~\ref{asumsiP} and let $\mathcal R_{P}(\phi)$ be defined by \eqref{eq:potential-remainder}.
Fix an integer $k\ge 0$ and let $\mathcal Z$ be any collection of \emph{first-order} admissible commutators (for instance, the families used in the Kerr and Schwarzschild parts of this paper).
Then there exists a constant $C=C(k,N_{P},\{C_{P,j}\}_{0\le j\le k})$ such that for every smooth $\phi$ and every hypersurface $\Sigma$ in the admissible foliation,
\begin{equation}\label{eq:potential-tame-L2}
\sum_{|\alpha|\le k}\|\mathcal Z^{\alpha}\mathcal R_{P}(\phi)\|_{L^{2}(\Sigma)}
\le
C\,(1+\|\phi\|_{L^{\infty}(\Sigma)}^{2N_{P}})\,\|\phi\|_{L^{\infty}(\Sigma)}^{2}\,
\sum_{|\alpha|\le k}\|\mathcal Z^{\alpha}\phi\|_{L^{2}(\Sigma)}.
\end{equation}
The same bound holds with $L^{2}(\Sigma)$ replaced by $L^{2}(\mathcal R(\tau_{0},\tau_{1}))$ and $L^{\infty}(\Sigma)$ replaced by $L^{\infty}(\mathcal R(\tau_{0},\tau_{1}))$.
\end{lemma}

\begin{definition}[Massive Maxwell-Higgs energy and bulk norms on slowly rotating Kerr]\label{def:slow-massive-energy-norms}
Let \(m^2>0\) and let
\begin{equation}
 P(\phi,\bar\phi)=m^2|\phi|^2+P_{\ge p}(|\phi|^2),
 \qquad
 P_{\ge p}(s)=\sum_{n=p}^{N_P}\alpha_n s^n,
 \qquad p\ge2,
 \qquad \alpha_n\ge0.
\end{equation}
Let \(Y_{\rm red}\) be a smooth redshift vector field supported in a fixed neighbourhood of the future horizon, transversal to \(\mathcal H^+\), and chosen so that the redshift bulk is positive.  We use the horizon-adapted timelike multiplier
\begin{equation}\label{eq:Tchi-slow-definition}
 T_\chi=T+\chi(r)\Omega_H\Phi+Y_{\rm red},
 \qquad
 T=\partial_t,
 \qquad
 \Phi=\partial_\varphi,
 \qquad
 \Omega_H=\frac{a}{r_+^2+a^2},
\end{equation}
where \(\chi\) equals one near the future horizon and vanishes away from a fixed compact radial set, while \(Y_{\rm red}\) is zero in the asymptotic region.  Thus \(T_\chi=T\) for large \(r\).  For \(|a|/M\) sufficiently small this vector field is uniformly future timelike.  The added redshift component is essential: the horizon generator \(T+\Omega_H\Phi\) is null on \(\mathcal H^+\) and by itself would not control the horizon-transversal derivative.  Let \((L,n,m,\bar m)\) be a redshift-regular Carter null frame and set
\begin{equation}\label{eq:slow-maxwell-scalars}
 \Phi_{+1}=F(L,m),
 \qquad
 \Phi_{-1}=F(\bar m,n),
 \qquad
 \Phi_0=\frac12\{F(L,n)+F(\bar m,m)\},
 \qquad
 \Upsilon=q^2\Phi_0,
\end{equation}
with \(q=r-ia\cos\theta\).  For a smooth Lorenz-gauge zero-sector solution, normalized in the residual Lorenz class as in Lemma~\ref{lem:slow-lorenz-potential-control}, define first the componentwise Lorenz-potential energy
\begin{equation}\label{eq:slow-lorenz-potential-energy}
E_N^{A}(\tau):=
\sum_{\nu=0}^{3}\sum_{|I|\le N}\int_{\Sigma_\tau}
\left(|\nabla Z^IA_\nu|_N^2+r^{-2}|Z^IA_\nu|^2\right)\,\dd\mu_{\Sigma_\tau}.
\end{equation}
Then set
\begin{equation}
E_{N}^{MH,m}(\tau)
:=
\sum_{|I|\le N}\int_{\Sigma_\tau}
\Bigl(
 |\mathcal L_{Z}^{I}F|_{N}^{2}
 +|D Z^{I}\phi|_{N}^{2}
 +m^{2}|Z^{I}\phi|^{2}
\Bigr)\,\dd\mu_{\Sigma_\tau}
 +\int_{\Sigma_\tau}P_{\ge p}(|\phi|^{2})\,\dd\mu_{\Sigma_\tau}.
\label{eq:slow-energy-massive-def}
\end{equation}
The norm \(|\cdot|_N\) is the nondegenerate frame norm associated with the redshift frame.  The potential term is uncommuted; the commuted nonlinear potential force is estimated in the source norm by Lemma~\ref{lem:potential-tame}.  The Fackerell-Ipser energy is
\begin{equation}\label{eq:slow-FI-energy-def}
 E_N^{FI}(\tau):=
 \sum_{|I|\le N-2}\int_{\Sigma_\tau}
 \Bigl(|\nabla Z^I\Upsilon|_N^2+r^{-2}|Z^I\Upsilon|^2\Bigr)
 \dd\mu_{\Sigma_\tau}.
\end{equation}
For \(0\le \tau\le T\) the spacetime norms are the sums, over \(|I|\le N-2\), of the following quantities.  With
\begin{equation}
 w_r=\frac{\Delta}{(r^2+a^2)^2},
\end{equation}
we first define the auxiliary Lorenz-potential bulk by
\begin{equation}\label{eq:slow-A-bulk-def}
 \mathcal B_N^A(0,T)
 :=
 \sum_{\nu=0}^{3}\sum_{|I|\le N}
 \|Z^I A_\nu\|_{\mathcal X^{\rm wave}_{0}(0,T)}^2,
\end{equation}
where \(\mathcal X^{\rm wave}_{0}\) is the massless scalar redshift-Morawetz-far-field solution norm used in Definition~\ref{def:lin-estimates}.  This quantity is not a gauge-invariant observable; it is attached to the normalized Lorenz representative and is used only to control the reduced scalar source.  We then set
\begin{eqnarray}
B_{\pm}^{MH}[F,\phi]
&:=&
\int_{0}^{T}\!\!\int_{\Sigma_\tau}
 w_r\Bigl(|\Phi_{+1}|^2+|\Phi_{-1}|^2+|D_r\phi|^2
 +r^{-2}|\nabla_{\mathbb S^2}\phi|^2\Bigr)\,\dd\mu,
\label{eq:slow-Bpm-def}\\
B_{0,m}^{MH}[\Upsilon,\phi]
&:=&
\int_{0}^{T}\!\!\int_{\Sigma_\tau}
\left[
 \frac{M}{r^4}|\Upsilon|^2
 +\frac1{r^2}\left(|D_t\phi|^2+m^2|\phi|^2+P_{\ge p}(|\phi|^2)\right)
\right]\dd\mu,
\label{eq:slow-B0-def}\\
B_{2,0,m}^{MH}[\Upsilon,\phi]
&:=&
\int_{0}^{T}\!\!\int_{\Sigma_\tau}
\left[
 \frac{M\Delta}{(r^2+a^2)^2}\left(|\partial_r\Upsilon|^2
 +|\nabla_{\mathbb S^2}\Upsilon|^2\right)
 +\frac1r|\nabla_{\mathbb S^2}\phi|^2
\right]\dd\mu,
\label{eq:slow-B20-def}\\
B_1^{MH}[\Upsilon,\phi]
&:=&
\int_{r_+}^{\infty}\chi_{tr}(r)
\left|\int_0^T\!\!\int_{\mathbb S^2}
\left\{
 \Im(\overline\Upsilon\,\partial_t\Upsilon)
 +\Im(\overline\phi\,D_t\phi)
\right\}\dd\omega\dd t\right|\dd r,
\label{eq:slow-B1-def}
\end{eqnarray}
where \(\chi_{tr}\) is supported in a small fixed neighborhood of the photon sphere.  Finally,
\begin{equation}\label{eq:slow-bulk-massive-def}
\begin{aligned}
 \mathcal B_{N}^{MH,m}(0,T)
 :=
 \mathcal B_N^A(0,T)+
 \sum_{|I|\le N-2}
 \Bigl(&
 B_{\pm}^{MH}[\mathcal L_Z^IF,Z^I\phi]
 +B_{0,m}^{MH}[Z^I\Upsilon,Z^I\phi]
 +B_{1}^{MH}[Z^I\Upsilon,Z^I\phi]\\
 &+B_{2,0,m}^{MH}[Z^I\Upsilon,Z^I\phi]
 \Bigr).
\end{aligned}
\end{equation}
Changing the redshift frame or the cutoffs within the admissible class changes these norms by constants independent of \(T\).
\end{definition}

\begin{lemma}[Massive Sobolev control in the energy bootstrap]\label{lem:slow-massive-sobolev-bootstrap}
Let \(N\ge6\), and let \((A,F,\phi)\) be a smooth zero-sector solution on \(0\le\tau\le T\).  If
\begin{equation}
 \mathfrak E_N^{(m)}(T):=
 \sup_{0\le s\le T}\bigl(E_N^{MH,m}(s)+E_N^{FI}(s)\bigr)
 +\mathcal B_N^{MH,m}(0,T)
 \le R^2,
\end{equation}
then the lower-order fields satisfy
\begin{equation}\label{eq:slow-pointwise-bootstrap-generic}
 \sum_{|J|\le N-3}
 \left(
 \|Z^J\phi\|_{L^\infty(\Sigma_\tau)}
 +\|D Z^J\phi\|_{L^\infty(\Sigma_\tau)}
 +\|\mathcal L_Z^JF\|_{L^\infty(\Sigma_\tau)}
 \right)
 \le C_{N,M}R
\end{equation}
for every \(0\le\tau\le T\).  The same bound holds locally for the normalized Lorenz potential after applying Lemma~\ref{lem:slow-lorenz-potential-control}.  In particular, after increasing \(C_b\) by a fixed constant, the bootstrap \(\mathfrak E_N^{(m)}(T)\le4C_b^2\varepsilon^2\) implies
\begin{equation}\label{eq:slow-pointwise-bootstrap}
\begin{split}
 \sum_{|J|\le N-3}
 \left(
 \|Z^JA\|_{W^{1,\infty}(\Sigma_\tau)}
 +\|\mathcal L_Z^JF\|_{L^\infty(\Sigma_\tau)}
 +\|Z^J\phi\|_{L^\infty(\Sigma_\tau)}
 +\|D Z^J\phi\|_{L^\infty(\Sigma_\tau)}
 \right)
 &\le C_b\varepsilon,\\
 &\hspace{-1em}0\le\tau\le T.
\end{split}
\end{equation}
\end{lemma}

\begin{lemma}[Lorenz-potential control from the curvature]\label{lem:slow-lorenz-potential-control}
In the zero electric sector fix the residual Lorenz freedom by requiring the boundary trace of the gauge parameter to vanish and by imposing the slice normalization that removes the harmonic part of the spatial potential.  Then, for every \(N\ge6\), the Lorenz representative satisfies on each Kerr-star slice
\begin{equation}\label{eq:slow-potential-control}
 \|A\|_{H^{N+1}(\Sigma_\tau)}+
 \|\nabla_nA\|_{H^{N}(\Sigma_\tau)}
 \le C\left(
 \sum_{|I|\le N}\|\mathcal L_Z^IF\|_{L^2(\Sigma_\tau)}
 +\|J\|_{H^{N-1}(\Sigma_\tau)}
 \right),
\end{equation}
where \(J_\mu=2\Im(\bar\phi D_\mu\phi)\).  In the small-data regime the last term is quadratic in the massive energy norm and therefore perturbative.  The estimate is used only to continue the Lorenz representative; the energy identity itself is formulated in the gauge-invariant variables \((F,\phi)\).
\end{lemma}

\begin{equation}\label{eq:slow-fi-forced}
 (\square_{g_{M,a}}-V_{FI})\Upsilon=N_{MH},
 \qquad
 V_{FI}=-2Mq^{-3}.
\end{equation}

\begin{lemma}[Perturbative Fackerell-Ipser source bound]\label{lem:slow-FI-source-bound}
Assume on \([0,T]\) the bootstrap \(\mathfrak E_N^{(m)}(T)\le4C_b^2\varepsilon^2\).  Let \(W(r)\) be the dual Fackerell-Ipser Morawetz weight associated with \(B^{MH}_{2,0,m}\).  Then
\begin{equation}\label{eq:slow-FI-source-bound}
 \sum_{|I|\le N-2}\int_{\Omega(0,T)} W(r)|Z^IN_{MH}|^2\,\dd\mu
 \le C(C_b\varepsilon)^2\mathcal B_N^{MH,m}(0,T).
\end{equation}
The scalar source generated by \(\mathcal R_P\) obeys the corresponding bound in the massive scalar source norm \(\mathcal S_N^{(m)}(0,T)\) appearing in \(\ME_N^{(m)}(M,a)\).
\end{lemma}

\begin{proposition}[Core estimates on a finite slab under \(\ME^{(m)}_N\)]\label{prop:slow-core-estimates-full}
Let \((A,F,\phi)\) be a smooth Lorenz-gauge zero-sector Maxwell-Higgs solution on \(0\le\tau\le T\), and assume \(\ME^{(m)}_N(M,a)\).  Under the pointwise bootstrap supplied by Lemma~\ref{lem:slow-massive-sobolev-bootstrap}, the following estimates hold for \(|a|/M\) sufficiently small:
\begin{equation}\label{eq:slow-core-energy}
E^{MH,m}(T)
\le
C E^{MH,m}(0)+C\frac{|a|}{M}
\Bigl(B_{\pm}^{MH}+B_{0,m}^{MH}\Bigr),
\end{equation}
\begin{equation}\label{eq:slow-core-morawetz}
B_{\pm}^{MH}
\le
C\Bigl(E^{MH,m}(T)+E^{MH,m}(0)+B_{0,m}^{MH}\Bigr),
\end{equation}
\begin{eqnarray}
E^{FI}(T)
&\le&
C\Bigl(E^{FI}(0)+E^{MH,m}(0)
+\frac{|a|}{M}(B_{\pm}^{MH}+B_{0,m}^{MH}+B_1^{MH}+B_{2,0,m}^{MH})\Bigr)
\nonumber\\
&&+C(C_b\varepsilon)^2\mathcal B_N^{MH,m}(0,T),
\label{eq:slow-core-fi}\\
B_{0,m}^{MH}+B_{2,0,m}^{MH}
&\le&
C\Bigl(E^{FI}(T)+E^{FI}(0)+E^{MH,m}(T)+E^{MH,m}(0)
+\frac{|a|}{M}B_{\pm}^{MH}\Bigr)
\nonumber\\
&&+C(C_b\varepsilon)^2\mathcal B_N^{MH,m}(0,T),
\label{eq:slow-core-fi-mor}\\
B_1^{MH}
&\le&
C\Bigl(E^{FI}(T)+E^{FI}(0)+E^{MH,m}(T)+E^{MH,m}(0)
+B_{0,m}^{MH}+B_{2,0,m}^{MH}\Bigr)
\nonumber\\
&&+C\frac{|a|}{M}\mathcal B_N^{MH,m}(0,T)
+C(C_b\varepsilon)^2\mathcal B_N^{MH,m}(0,T).
\label{eq:slow-core-trap}
\end{eqnarray}
The same bounds hold after summing over \(|I|\le N-2\), with \(E^{MH,m}\), \(E^{FI}\), and the \(B\)-terms replaced by their commuted versions.
\end{proposition}

\begin{proof}
The proof is by the vector-field identities in the redshift frame, with the massive scalar Morawetz/boundedness component supplied by \(\ME^{(m)}_N(M,a)\).  The computations in this section control the Maxwell, Fackerell-Ipser, source, and perturbative coupling parts and show that the nonlinear terms fit the source norm of the condition.

For \eqref{eq:slow-core-energy}, integrate the divergence identity for \(T^{MH}_{\alpha\beta}T_\chi^\beta\).  The Killing fields \(T\) and \(\Phi\) have zero deformation tensor.  The redshift part \(Y_{\rm red}\) contributes a positive near-horizon bulk, while the only indefinite terms come from the compact cutoffs and are proportional to \(\Omega_H=O(a/M^2)\).  In the redshift frame the indefinite part satisfies
\begin{equation}
\left|T^{MH}_{\alpha\beta}\nabla^{(\alpha}T_\chi^{\beta)}\right|_{\rm ind}
\le
C\frac{|a|}{M}w_r
\Bigl(|\Phi_{+1}|^2+|\Phi_0|^2+|\Phi_{-1}|^2+|D\phi|^2+m^2|\phi|^2+P_{\ge p}\Bigr).
\end{equation}
The positive redshift bulk is kept on the left or discarded, the displayed error is bounded by \((|a|/M)(B_{\pm}^{MH}+B_{0,m}^{MH})\), and the boundary fluxes are equivalent to \(E^{MH,m}\).  This proves \eqref{eq:slow-core-energy}.

For \eqref{eq:slow-core-morawetz}, take the radial multiplier
\begin{equation}
 \widehat A=f(r)\partial_r,
 \qquad
 f(r)=-\frac{\Delta}{r^2+a^2}\left(1-\frac{3M}{r}\right).
\end{equation}
Use the conformal form of the bulk identity with \(\Omega^{-2}=\Sigma\Delta/(r^2+a^2)^2\):
\begin{equation}
\operatorname{Bulk}_{\widehat A}^{MH}
=-\frac12\int T^{MH}_{\alpha\beta}\mathcal L_{\widehat A}(\Omega^{-2}g^{\alpha\beta})\,\dd\mu
+\frac12\int \Omega^{-2}T^{MH\,\gamma}{}_{\gamma}\widehat A(\Omega^2)\,\dd\mu.
\end{equation}
A direct Kerr coefficient calculation gives
\begin{eqnarray}
\mathcal L_{\widehat A}(\Omega^{-2}g^{\alpha\beta})
&=&
\left(f\partial_r\frac{\Delta^2}{(r^2+a^2)^2}
 -2\frac{\Delta^2}{(r^2+a^2)^2}f'\right)\partial_r^\alpha\partial_r^\beta
\nonumber\\
&&
+f\partial_rV_L(\Theta^\alpha\Theta^\beta+\Phi_{PNV}^\alpha\Phi_{PNV}^\beta)
+O\left(\frac{|a|}{M}w_r\right),
\end{eqnarray}
where \(V_L\) is the angular Kerr potential and the last term denotes the \(t\varphi\) and \(\varphi\varphi\) errors.  The radial coefficient is positive, comparable to \(\Delta^3/(r^2+a^2)^3\), and the angular coefficient is positive away from \(|r-3M|\lesssim r_1\), up to a controlled \(Mr^{-4}|\Upsilon|^2\) loss.  The trace term is bounded by \(B_{0,m}^{MH}\), since
\begin{equation}
 T^{MH\,\gamma}{}_{\gamma}=-2|D\phi|^2-4P(\phi,\bar\phi).
\end{equation}
Thus
\begin{equation}
 B_{\pm}^{MH}\le C\operatorname{Bulk}_{\widehat A}^{MH}+C B_{0,m}^{MH}.
\end{equation}
The boundary terms of the \(\widehat A\)-current are bounded by \(E^{MH,m}(0)+E^{MH,m}(T)\), which gives \eqref{eq:slow-core-morawetz}.

For \eqref{eq:slow-core-fi}, use the stress tensor of the complex scalar \(\Upsilon\) solving \eqref{eq:slow-fi-forced}:
\begin{equation}
T^{FI}_{\alpha\beta}[\Upsilon]
=\Re(\nabla_\alpha\Upsilon\overline{\nabla_\beta\Upsilon})
-\frac12g_{\alpha\beta}\left(|\nabla\Upsilon|^2+\Re(V_{FI})|\Upsilon|^2\right).
\end{equation}
Its divergence is
\begin{eqnarray}
\nabla^\alpha T^{FI}_{\alpha\beta}
&=&
\Re(N_{MH}\overline{\nabla_\beta\Upsilon})
+\frac12\partial_\beta\Re(V_{FI})|\Upsilon|^2
\nonumber\\
&&+\Im(V_{FI})\Im(\Upsilon\overline{\nabla_\beta\Upsilon}).
\end{eqnarray}
Contracting with \(T_\chi\), integrating, and using \(|T_\chi\Re V_{FI}|+|\Im V_{FI}|\le C|a|Mr^{-4}\), one obtains the endpoint FI energy bound with \((|a|/M)(B_{0,m}^{MH}+B_1^{MH}+B_{2,0,m}^{MH})\) errors.  The source term satisfies, for every \(\delta>0\),
\begin{equation}
\left|\int \Re(N_{MH}\overline{T_\chi\Upsilon})\,\dd\mu\right|
\le \delta B_{2,0,m}^{MH}
+C_\delta(C_b\varepsilon)^2\mathcal B_N^{MH,m}(0,T)
\end{equation}
by Lemma~\ref{lem:slow-FI-source-bound}; choosing \(\delta\) small gives \eqref{eq:slow-core-fi}.

For \eqref{eq:slow-core-fi-mor}, apply the radial FI current
\begin{equation}
\mathcal J^\alpha_{X,q_0}
=T^{FI\,\alpha}{}_{\beta}X^\beta
+\frac12q_0\nabla^\alpha(|\Upsilon|^2)
-\frac12(\nabla^\alpha q_0)|\Upsilon|^2,
\qquad X=\widehat A,
\end{equation}
with the zeroth-order correction \(q_0(r)\).  The direct divergence formula is
\begin{eqnarray}
\nabla_\alpha\mathcal J^\alpha_{X,q_0}
&=&T^{FI}_{\alpha\beta}\nabla^{(\alpha}X^{\beta)}
+\frac12q_0|\nabla\Upsilon|^2
-\frac14(\square_gq_0)|\Upsilon|^2
-\frac12X(\Re V_{FI})|\Upsilon|^2
\nonumber\\
&&+\Im(V_{FI})\Im(\overline\Upsilon X\Upsilon)
+\Re(N_{MH}\overline{X\Upsilon})
+q_0\Re(N_{MH}\overline\Upsilon).
\end{eqnarray}
The source-free part is the slowly rotating Fackerell-Ipser Morawetz density and controls \(B_{0,m}^{MH}+B_{2,0,m}^{MH}\), up to endpoint FI and Maxwell-Higgs energies and an \((|a|/M)B_{\pm}^{MH}\) error.  The two source terms are controlled by the dual weight \(W\):
\begin{equation}
 |I_{src}|
\le \delta(B_{0,m}^{MH}+B_{2,0,m}^{MH})
+C_\delta\sum_{|I|\le N-2}\int W|Z^IN_{MH}|^2\dd\mu.
\end{equation}
Lemma~\ref{lem:slow-FI-source-bound} and smallness of \(C_b\varepsilon\) give \eqref{eq:slow-core-fi-mor}.

For \eqref{eq:slow-core-trap}, insert the cutoff \(\chi_{tr}\) in \eqref{eq:slow-B1-def}.  Cauchy-Schwarz on \(\mathbb S^2\), followed by Young's inequality, gives
\begin{equation}
B_1^{MH}\le C B_{0,m}^{MH}
+C\eta\int_{\Omega(0,T)}\chi_{tr}|\partial_t\Upsilon|^2\dd\mu.
\end{equation}
The local elliptic estimate for \((\square_g-V_{FI})\Upsilon=N_{MH}\) on the compact radial slab \(|r-3M|\le2r_1\) yields
\begin{eqnarray}
\int\chi_{tr}|\partial_t\Upsilon|^2\dd\mu
&\le&
C\{E^{FI}(0)+E^{FI}(T)+B_{0,m}^{MH}+B_{2,0,m}^{MH}\}
\nonumber\\
&&+C\frac{|a|}{M}(B_{\pm}^{MH}+B_{0,m}^{MH}+B_{2,0,m}^{MH})
+C\sum_{|I|\le N-2}\int W|Z^IN_{MH}|^2\dd\mu.
\end{eqnarray}
Substituting Lemma~\ref{lem:slow-FI-source-bound}, and then choosing \(\eta\) and \(C_b\varepsilon\) small in the global bootstrap, gives \eqref{eq:slow-core-trap}.  All commuted estimates follow from the same identities applied to \(Z^IF\), \(Z^I\phi\), and \(Z^I\Upsilon\); commutators with Kerr coefficients have one lower differential order and one extra radial decay, and therefore are absorbed by the displayed lower-order bulk terms.
\end{proof}

\end{document}
